% 题干转写：完成
% 解析转写：完成
% 图形核对：完成
\section{第36届复赛第4题：微孔竹板亥姆霍兹共振吸声}
\subsection{题干}
微孔竹板上的周期性长方形微孔可使孔中空气柱发生简谐振动；单个微孔及其后方空气层可简化为亥姆霍兹共振器，如图\ref{fig:cpho36q04}。微孔深度为 $l$，孔横截面积为 $S$，单孔后腔体体积为 $V_0$。声波传播视为绝热过程，满足 $v_s=\sqrt{\kappa/\rho}$，其中 $\kappa=-V\Delta p/\Delta V$。已知标准状态 $T=273\,\mathrm{K}$、$p=1.01\times10^5\,\mathrm{Pa}$，空气摩尔质量 $M_{\rm mol}=29.0\,\mathrm{g/mol}$、热容比 $\gamma=7/5$、气体常量 $R=8.31\,\mathrm{J/(K\cdot mol)}$。求：（1）标准状态下空气的密度和声速；（2）共振器吸收声波的频率；（3）吸收 $120\,\mathrm{Hz}$ 和 $200\,\mathrm{Hz}$ 时两种微孔面积之比。
\begin{figure}[htbp]\centering
\begin{tikzpicture}[scale=0.83,>=Stealth,font=\small]
  % The left panel preserves the periodic perforation feature without reproducing the source photograph.
  \draw[fill=brown!22,thick] (0,0) rectangle (3.15,2.45);
  \foreach \x in {0.22,0.48,...,2.92} {
    \foreach \y in {0.22,0.48,...,2.22} {
      \fill[black!65] (\x,\y) circle (0.026);
    }
  }
  \node[align=center,below] at (1.58,-0.10) {图a：周期性微孔竹板\\（结构示意）};

  % Cross section of the panel and its rear air layer.
  \begin{scope}[xshift=4.0cm]
    \draw[fill=gray!20,thick] (0,0) rectangle (5.4,0.30);
    \draw[fill=gray!12,thick] (0,0.30) rectangle (5.4,1.78);
    \foreach \x in {0.18,0.55,...,5.15} {
      \draw[fill=gray!60,draw=black] (\x,1.78) rectangle +(0.24,0.98);
    }
    \draw[thick] (0,2.76) -- (5.4,2.76) -- (5.4,0) -- (0,0) -- cycle;
    \draw[densely dashed] (1.82,0.05) -- (1.82,2.73);
    \draw[densely dashed] (2.34,0.05) -- (2.34,2.73);
    \draw[<->] (1.82,3.12) -- (2.34,3.12) node[midway,above] {$S$};
    \draw[<->] (5.78,1.78) -- (5.78,2.76) node[midway,right] {$l$};
    \node at (3.65,1.05) {空气层};
    \node[below] at (2.7,-0.10) {微孔板与后腔截面};

    % Corresponding lumped Helmholtz resonator.
    \begin{scope}[xshift=6.7cm]
      \draw[fill=gray!20,thick] (0,0) rectangle (1.62,1.42);
      \draw[fill=gray!55,thick] (0,0) rectangle (1.62,0.25);
      \draw[thick] (0.55,1.42) -- (0.55,2.38) -- (1.07,2.38) -- (1.07,1.42);
      \draw[<->] (1.32,1.42) -- (1.32,2.38) node[midway,right] {$l$};
      \draw[<->] (0.55,2.70) -- (1.07,2.70) node[midway,above] {$S$};
      \node at (0.81,0.80) {$V_0$};
      \node[below,align=center] at (0.81,-0.10) {亥姆霍兹\\共振器模型};
    \end{scope}
  \end{scope}
\end{tikzpicture}

\caption{第36届复赛第4题微孔竹板及亥姆霍兹共振器的 TikZ 重绘。}\label{fig:cpho36q04}
\end{figure}

\subsection{参考解答整理}
\paragraph{（1）标准状态下的密度与声速。} 理想气体状态方程为
\[
pV=\frac{M}{M_{\rm mol}}RT,
\qquad
\rho=\frac{M}{V}=\frac{M_{\rm mol}p}{RT}.
\]
故标准状态空气密度为
\[
\boxed{\rho_0=1.29\,\mathrm{kg/m^3}.}
\]
绝热过程满足 $pV^\gamma=\text{常量}$。微分后可得
\[
\kappa=-V\frac{\Delta p}{\Delta V}=\gamma p,
\qquad
v_s=\sqrt{\frac{\gamma p}{\rho}}.
\]
因此标准状态下
\[
\boxed{v_{s0}=\sqrt{\frac{(7/5)\times1.01\times10^5}{1.29}}=331\,\mathrm{m/s}.}
\]

\paragraph{（2）亥姆霍兹共振频率。} 微孔内空气柱质量为 $m=\rho Sl$。令空气柱向外位移为 $x$，则腔体体积改变量为 $\Delta V=Sx$，压强变化和回复力为
\[
\Delta p=-\frac{\gamma pSx}{V},
\qquad
F=S\Delta p=-\frac{\gamma pS^2}{V}x.
\]
由 $m\ddot{x}=F$ 得
\[
\ddot{x}=-\omega^2x,
\qquad
\omega=\sqrt{\frac{\gamma pS}{\rho lV}}.
\]
将 $V=V_0$ 与 $v_s=\sqrt{\gamma p/\rho}$ 代入，得到
\[
\boxed{f_0=\frac{v_s}{2\pi}\sqrt{\frac{S}{lV_0}}.}
\]

\paragraph{（3）两种孔面积之比。} 在 $l$、$V_0$ 和 $v_s$ 相同的条件下，$f\propto\sqrt S$，所以
\[
\frac{S_1}{S_2}=\left(\frac{f_{01}}{f_{02}}\right)^2
=\left(\frac{120}{200}\right)^2.
\]
故
\[
\boxed{\frac{S_1}{S_2}=0.36.}
\]

\subsection{转写依据}
题干、答案公式和图形分别依据 \texttt{CPHO36-Q04.md}、\texttt{CPHO36-Q04\_formula\_check\_v1.csv} 与题图核对入口。图a为照片型原图，成品以微孔周期结构示意替代，不嵌入原题截图。评分分布为第（1）问18分、第（2）问16分、第（3）问6分。
